\subsection{Problemática}

Cuando varias personas desarrollan simultáneamente un mismo videojuego, cada integrante tiende a resolver problemas equivalentes --- nombrar una variable, estructurar una clase, registrar un error, escribir una prueba --- de una manera distinta. En ausencia de un criterio compartido, el código de un mismo proyecto termina mezclando convenciones de nombrado incompatibles, niveles de anidamiento excesivos, manejo inconsistente de excepciones y bitácoras improvisadas con \texttt{Console.WriteLine}. Esta heterogeneidad incrementa el tiempo que cualquier integrante necesita para comprender un módulo escrito por alguien más, dificulta la revisión de código y facilita la introducción de errores al modificar lógica que no se entiende por completo.

El problema se agrava en un proyecto como un juego de mesa digital en 2D, donde la lógica de turnos, movimientos, validación de reglas y estado del tablero se reparte entre varias clases que deben interactuar de forma predecible. Sin reglas explícitas de nombrado, manejo de errores, complejidad y pruebas unitarias, el equipo no cuenta con un criterio objetivo para decidir si una porción de código está lista para integrarse al resto del proyecto.
